% eq_gravity_derived.tex — Gravity as Derived Consequence of Ax 1 + Ax 4 (Symmetric Scaling)
% Single source of truth for the gravitational refractive index, Symmetric Gravity
% derivation, and alpha-invariance under gravitational strain. \input this wherever
% the gravity content is formally restated.
%
% PROVENANCE: This file documents content that historically lived in eq_axiom_3.tex
% (Scheme B framing of "Axiom 3 = Gravity"). Per axiom homologation 2026-04-27 +
% Vol 1 Ch 1's canonical Scheme A: Axiom 3 is the Minimum Reflection Principle
% (eq_axiom_3.tex; renamed 2026-05-16 from "Effective Action Principle" — the new
% name follows the substrate-observable, the boundary reflection |Γ|², rather
% than the interior Lagrangian density). The gravitational refractive index n(r)
% and Symmetric Gravity scaling are derived consequences of Axiom 1 (LC network) +
% Axiom 4 (dielectric saturation applied symmetrically to mu and epsilon). Newton's
% constant G is a calibration constant emerging from Machian thermodynamic
% equilibrium of the LC lattice (per Vol 1 Ch 1:14-21 and Vol 1 Ch 1 Section 3.2).
%
\begin{resultbox}{Macroscopic Gravity (Derived from Ax 1 + Ax 4 Symmetric Scaling)}
Newton's constant sets the Machian boundary impedance:
\begin{equation}
    G = \frac{\hbar c}{7\xi \cdot m_e^2}
    \label{eq:gravity_newton}
\end{equation}
where $\xi = 4\pi(R_H / \ell_{node})\alpha^{-2} \approx 8.15 \times 10^{43}$ is the Machian hierarchy coupling.
The gravitational refractive index is
\begin{equation}
    n(r) = 1 + \frac{2GM}{rc^2}
    \label{eq:gravity_refraction}
\end{equation}
with Symmetric Gravity (Axiom 4 dielectric saturation applied to both inductive and capacitive constituents):
$\mu_{eff} = \mu_0 \cdot n(r)$,\;
$\varepsilon_{eff} = \varepsilon_0 \cdot n(r)$,\;
$Z = Z_0$ (invariant).
\end{resultbox}

\noindent\textbf{Derived Consequence 1: $\alpha$ Invariance.}
Under Symmetric Gravity, both constitutive parameters scale by the same factor $n \cdot S$
(combining gravitational refraction with Axiom 4 saturation). The fine-structure constant is therefore
\textit{exactly invariant} under gravitational strain:
\begin{equation}
    \alpha = \frac{e^2}{4\pi\,\varepsilon_{local}\,\hbar\,c_{local}}
           = \frac{e^2}{4\pi\,(\varepsilon_0 n S)\,\hbar\,(c_0/(nS))}
           = \frac{e^2}{4\pi\,\varepsilon_0\,\hbar\,c_0}
           = \alpha_0
    \label{eq:alpha_invariance}
\end{equation}
Multi-species clock comparisons at different gravitational potentials predict a
\textbf{null result} for $\Delta\alpha/\alpha$, consistent with all current experimental bounds.

\noindent\textbf{Derived Consequence 2: Temporal vs.\ Spatial Lattice Compression.}
The principal radial strain $\varepsilon_{11} = 7GM/(c^2 r)$ compresses the lattice asymmetrically.
The refractive index decomposes into:
\begin{align}
    n_{temporal} &= 1 + \tfrac{2}{7}\,\varepsilon_{11} && (\text{bulk/coordinate-time propagation index, slope 2}) \nonumber \\
    n_{spatial} &= 1 + \tfrac{9}{7}\,\varepsilon_{11} && (\text{controls matter-wave parallax, C11})
    \label{eq:lattice_decomposition}
\end{align}
% Notation cleanup 2026-05-17: both indices carry the "1 +" DC unit for parallelism.
% Both reduce to n = 1 in flat vacuum (eps_11 -> 0). The asymmetric prior form
% (n_spatial without "1 +") was shorthand for the strain-contribution term only,
% which produced confusion in driver-script implementations. The parity violation
% Delta_n = n_spatial - n_temporal = (9/7 - 2/7) * eps_11 = eps_11 is unchanged
% by the notation cleanup.
The temporal index $n_{temporal}$ (slope 2) is the \emph{bulk/coordinate-time}
propagation index --- what a signal traversing the gradient accumulates (Shapiro-class
integrated delay), $\approx 1/g_{00}$. The genuine \emph{local} clock rate / gravitational
redshift is a \emph{distinct} (slope 1) quantity, the proper tick of a clock sitting at $r$:
$\sqrt{g_{00}} = \sqrt{S} \approx 1 - GM/(c^2 r)$, giving redshift
$z \approx GM/(c^2 r)$. The two are bridged by $z = (n_{temporal} - 1)/2$: a propagating
signal picks up $2\times$ the local clock effect.
The spatial component governs the matter-wave (electron de~Broglie) parallax measured by
C11-MACH-ZEHNDER ($n_s - n_t = \varepsilon_{11}$) and frame dragging.
% W1 walk-back 2026-06-05 (Class-C internal-coherence; basis: research/2026-06-05_gravity-ppn-coherence-result.md):
% the (9/7) spatial index does NOT control light deflection. The canonical light-deflection
% derivation (Ch 2 \S double_deflection -> 4GM/bc^2) uses the (2/7) TRANSVERSE Cosserat-shear
% index n_perp = 1 + (2/7)chi_vol, not n_spatial. The (9/7) index IS load-bearing for the
% C11 electron matter-wave parallax (249.64 rad). KEEP-BOTH: (9/7) retained; only the
% "light deflection" attribution corrected.
% W2 walk-back 2026-06-05 (Class-C internal-coherence; Grant bulk-vs-local adjudication;
% basis: research/2026-06-05_gravity-sign-frequency-modulation-result.md):
% the "controls clock rate, redshift" annotation on n_temporal conflated TWO distinct
% temporal quantities. n_temporal (slope 2) = BULK/coordinate-time propagation index
% (integrated Shapiro delay, ~1/g00). The LOCAL clock rate / redshift is slope 1:
% sqrt(g00) = sqrt(S) ~ 1 - GM/rc^2, z = GM/rc^2 (the c_shear clock). Bridge:
% z = (n_temporal - 1)/2. RELABEL ONLY: n_temporal value (2/7 eps_11) and all
% derivations unchanged; only the "clock rate, redshift" label disambiguated from the
% slope-1 local clock. KEEP-BOTH.
